% Source: printed pp. 475-478; PDF pp. 484-487.
% Titles and publication names remain in their original languages.
\renewcommand{\bibname}{Tài liệu tham khảo}
\begin{thebibliography}{46}
\raggedright
\addcontentsline{toc}{chapter}{Tài liệu tham khảo}
\bibitem{ref1} Abramsky, S. (2012) Relational databases and Bell's Theorem.
Có tại \url{http://arxiv.org/abs/1208.6416}.
\bibitem{ref2} Atiyah, M. (1989) Topological quantum field theories.
\emph{Publications Mathématiques de l'IHÉS} 68(1), 175-186.
\bibitem{ref3} Axler, S. (1997) \emph{Linear Algebra Done Right}. Ấn bản thứ 2.
New York: Springer.
\bibitem{ref4} Awodey, S. (2010) \emph{Category Theory}. Ấn bản thứ 2.
Oxford: Oxford University Press.
\bibitem{ref5} Bralow, H. (1961) Possible principles underlying the transformation
of sensory messages. Trong \emph{Sensory Communication}, W. Rosenblaith biên tập,
217-234. Cambridge, MA: MIT Press.
\bibitem{ref6} Baez, J.C.; Dolan, J. (1995) Higher-dimensional algebra and
topological quantum field theory. \emph{Journal of Mathematical Physics} 36:
6073-6105.
\bibitem{ref7} Baez, J.C.; Fritz, T.; Leinster, T. (2011) A characterization
of entropy in terms of information loss. \emph{Entropy} 13(11): 1945-1957.
\bibitem{ref8} Baez, J.C.; Stay, M. (2011) Physics, topology, logic and
computation: a Rosetta Stone. Trong \emph{New Structures for Physics},
B. Coecke biên tập, 95Ð172. Lecture Notes in Physics 813.
Heidelberg: Springer.
\bibitem{ref9} Brown, R.; Porter, T. (2006) Category Theory: An abstract setting
for analogy and comparison. Trong \emph{What Is Category Theory?},
G. Sica biên tập, 257-274. Advanced Studies in Mathematics and Logic.
Monza Italy: Polimetrica.
\bibitem{ref10} Brown, R.; Porter, T. (2003) Category theory and higher
dimensional algebra: potential descriptive tools in neuroscience.
Trong \emph{Proceedings of the International Conference on Theoretical
Neurobiology}, tập 1, 80-92.
% Source: printed p. 476; PDF p. 485.
\bibitem{ref11} Barr, M.; Wells, C. (1990) \emph{Category Theory for Computing
Science}. New York: Prentice Hall.
\bibitem{ref12} Biggs, N.M. (2004) \emph{Discrete Mathematics}.
New York: Oxford University Press.
\bibitem{ref13} Diaconescu, R. (2008) \emph{Institution-Independent Model Theory}.
Boston: Birkhäuser.
\bibitem{ref14} Döring, A.; Isham, C.J. (2008) A topos foundation for theories
of physics. I. Formal languages for physics. \emph{Journal of Mathematical
Physics} 49(5): 053515.
\bibitem{ref15} Ehresmann, A.C.; Vanbremeersch, J-P. (2007)
\emph{Memory Evolutive Systems: Hierarchy, Emergence, Cognition}.
Amsterdam: Elsevier.
\bibitem{ref16} Everett III, H. (1973) The theory of the universal wave function.
Trong \emph{The Many-Worlds Interpretation of Quantum Mechanics},
B.S. DeWitt và N. Graham biên tập, 3-140. Princeton, NJ:
Princeton University Press.
\bibitem{ref17} Goguen, J. (1992) Sheaf semantics for concurrent interacting
objects. \emph{Mathematical Structures in Computer Science} 2(2): 159-191.
\bibitem{ref18} Grothendieck, A.; Raynaud, M. (1971)
\emph{Revêtements étales et groupe fondamental. Séminaire de Géométrie
Algébrique du Bois Marie, 1960/61 (SGA 1)}. Lecture Notes in Mathematics 224.
Tiếng Pháp. New York: Springer.
\bibitem{ref19} Krömer, R. (2007) \emph{Tool and Object: A History and Philosophy
of Category Theory}. Boston: Birkhäuser.
\bibitem{ref20} Lambek, J. (1980) From $\lambda$-calculus to Cartesian closed
categories. Trong \emph{To H.B. Curry: Essays on Combinatory Logic, Lambda
Calculus and Formalism}, J.P. Seldin và J. Hindley biên tập, 376-402.
London: Academic Press.
\bibitem{ref21} Khovanov, M. (2000) A categorificiation of the Jones polynomial.
\emph{Duke Mathematical Journal} 101(3): 359-426.
\bibitem{ref22} Landry, E.; Marquis, J.-P. (2005) Categories in contexts:
Historical, foundational, and philosophical. \emph{Philosophia Mathematica}
13(1): 1-43.
\bibitem{ref23} Lawvere, F.W. (2005) An elementary theory of the category of sets
(long version) with commentary. \emph{Reprints in Theory and Applications of
Categories}, số 11, 1-35. Mở rộng từ \emph{Procedings of the National Academy
of Sciences} 1964; 52(6): 1506-1511.
\bibitem{ref24} Lawvere, F.W.; Schanuel, S.H. (2009)
\emph{Conceptual Mathematics. A First Introduction to Categories}.
Ấn bản thứ 2. Cambridge: Cambridge University Press.
% Source: printed p. 477; PDF p. 486.
\bibitem{ref25} Leinster, T. (2004) \emph{Higher Operads, Higher Categories}.
London Mathematical Society Lecture Note Series 298.
New York: Cambridge University Press.
\bibitem{ref26} Leinster, T. (2012) Rethinking set theory.
Có tại \url{http://arxiv.org/abs/1212.6543}.
\bibitem{ref27} Linsker, R. (1988) Self-organization in a perceptual network.
\emph{Computer} 21(3): 105-117.
\bibitem{ref28} MacKay, D.J. (2003)
\emph{Information Theory, Inference and Learning Algorithms}.
Cambridge: Cambridge University Press.
\bibitem{ref29} Mac Lane, S. (1998)
\emph{Categories for the Working Mathematician}. Ấn bản thứ 2.
New York: Springer.
\bibitem{ref30} Marquis, J.-P. (2009) \emph{From a Geometrical Point of View:
A Study in the History and Philosophy of Category Theory}.
New York: Springer.
\bibitem{ref31} Marquis, J.-P. (2013) Category theory.
Trong \emph{Stanford Encyclopedia of Philosophy} (ấn bản mùa hè),
E.N. Zalta biên tập. Có tại
\url{http://plato.stanford.edu/archives/spr2011/entries/category-theory}.
\bibitem{ref32} Minsky, M. (1985) \emph{The Society of Mind}.
New York: Simon and Schuster.
\bibitem{ref33} Moggi, E. (1991) Notions of computation and monads.
\emph{Information and Computation} 93(1): 52-92.
\bibitem{ref34} nLab. \url{http://ncatlab.org/nlab/show/HomePage}.
\bibitem{ref35} Penrose, R. (2005) \emph{The Road to Reality}.
New York: Knopf.
\bibitem{ref36} Radul, A.; Sussman, G.J. (2009) \emph{The Art of the Propagator}.
MIT Computer Science and Artificial Intelligence Laboratory Technical Report.
\bibitem{ref37} Simmons, H. (2011) \emph{An Introduction to Category Theory}.
New York: Cambridge University Press.
\bibitem{ref38} Spivak, D.I. (2012) Functorial data migration.
\emph{Information and Computation} 217 (tháng Tám): 31-51.
\bibitem{ref39} Spivak, D.I. (2013) Database queries and constraints via lifting
problems. \emph{Mathematical structures in computer science}, 1-55.
Có tại \url{http://arxiv.org/abs/1202.2591}.
\bibitem{ref40} Spivak, D.I. (2012) Kleisli database instances.
Có tại \url{http://arxiv.org/abs/1209.1011}.
% Source: printed p. 478; PDF p. 487.
\bibitem{ref41} Spivak, D.I. (2013) The operad of wiring diagrams: Formalizing
a graphical language for databases, recursion, and plug-and-play circuits.
Có tại \url{http://arxiv.org/abs/1305.0297}.
\bibitem{ref42} Spivak, D.I.; Giesa, T.; Wood, E.; Buehler, M.J. (2011)
Category-theoretic analysis of hierarchical protein materials and social networks.
\emph{PLoS ONE} 6(9): e23911.
\bibitem{ref43} Spivak, D.I.; Kent, R.E. (2012)
Ologs: A categorical framework for knowledge representation.
\emph{PLoS ONE} 7(1): e24274.
\bibitem{ref44} Weinberger, S. (2011) What is ... persistent homology?
\emph{Notices of the AMS} 58(1): 36-39.
\bibitem{ref45} Weinstein, A. (1996) Groupoids: Unifying internal and external
symmetry. \emph{Notices of the AMS} 43(7): 744-752.
\bibitem{ref46} Wikipedia. Truy cập trong khoảng từ ngày 6 tháng 12 năm 2012
đến ngày 31 tháng 12 năm 2013.
\end{thebibliography}
